% ============================================================================
% APPENDIX DC: LIT-CARD - Labor Economics & Natural Experiments
% ============================================================================
% Category: Literature Integration (LIT-)
% Language: English
% EBF Integration: Causal inference, labor markets, and policy evaluation
% ============================================================================

\chapter*{Appendix DC: David Card Research Integration}
\addcontentsline{toc}{chapter}{Appendix DC: LIT-CARD - Labor Economics \& Natural Experiments}
\label{app:card}

% ============================================================================
% HEADER BLOCK
% ============================================================================
\begin{tcolorbox}[colback=blue!5!white, colframe=blue!75!black, title=Appendix Metadata]
\textbf{Code:} DC \\
\textbf{Category:} LIT- (Literature Integration) \\
\textbf{Title:} Labor Economics \& Natural Experiments \\
\textbf{Version:} 1.0 \\
\textbf{Status:} Complete \\
\textbf{Primary Author Focus:} David Card (Nobel Prize 2021) \\
\textbf{Papers Covered:} 30 \\
\textbf{Dependencies:} DOMAIN-LABOR, METHOD-EXPERIMENT, LIT-HECKMAN (N), DOMAIN-EDUCATION \\
\textbf{EBF Connections:} Causal identification, labor market parameters, policy evaluation
\end{tcolorbox}

% ============================================================================
% CROSS-REFERENCE MAP
% ============================================================================
\begin{tcolorbox}[colback=green!5!white, colframe=green!75!black, title=Cross-Reference Map]
\textbf{Outgoing References:}
\begin{itemize}[nosep]
    \item DOMAIN-LABOR: Labor market theory and evidence
    \item DOMAIN-EDUCATION: Returns to schooling
    \item METHOD-EXPERIMENT: Natural experiment methodology
    \item LIT-HECKMAN (N): Selection and causal inference
    \item LIT-DELLAVIGNA (DV): Field experiment collaboration
    \item CONTEXT-DEMOGRAPHIC: Immigration and gender effects
\end{itemize}

\textbf{Incoming References:}
\begin{itemize}[nosep]
    \item Chapter 4 (Methodology): Design-based research
    \item Chapter 14 (Applications): Labor policy design
    \item METHOD-META: Active labor market meta-analyses
\end{itemize}
\end{tcolorbox}

% ============================================================================
% ABSTRACT
% ============================================================================
\section*{Abstract}

David Card revolutionized empirical economics through his development of \textbf{natural experiment} methodology and its application to fundamental labor market questions. His 2021 Nobel Prize recognized ``empirical contributions to labour economics.'' This appendix integrates Card's 30 key papers into the EBF framework, with particular emphasis on:
\begin{enumerate}
    \item \textbf{Minimum wage effects}: The Card-Krueger revolution
    \item \textbf{Immigration and labor markets}: The Mariel Boatlift natural experiment
    \item \textbf{Returns to education}: Instrumental variable approaches
    \item \textbf{Wage inequality}: Firm effects and labor market structure
    \item \textbf{Active labor market policies}: Meta-analytic evidence
\end{enumerate}

% ============================================================================
% QUICK REFERENCE
% ============================================================================
\begin{tcolorbox}[colback=yellow!10!white, colframe=orange!75!black, title=Quick Reference: Key Card Parameters]
\begin{tabular}{lll}
\textbf{Parameter} & \textbf{Value} & \textbf{Source} \\
\hline
Minimum wage employment effect & $\approx 0$ & Card \& Krueger (1994) \\
Immigration wage elasticity & -0.1 to -0.2 & Card (1990, 2001) \\
Returns to education & 8-12\% per year & Card (1999, 2001) \\
Firm wage premium share & 20-25\% of variance & Card et al. (2013) \\
Peer salary satisfaction effect & -15\% if underpaid & Card et al. (2012) \\
\end{tabular}
\end{tcolorbox}

% ============================================================================
% FUNDAMENTAL QUESTION
% ============================================================================
% -----------------------------------------------------------------------------
% APPENDIX SCOPE BOX
% -----------------------------------------------------------------------------

\begin{tcolorbox}[
    colback=orange!5!white,
    colframe=orange!75!black,
    title={\textbf{Appendix Scope: Ziel / In-Scope / Out-of-Scope / Constraints}},
    fonttitle=\bfseries
]

\textbf{Ziel (Objective):}
\begin{itemize}[nosep]
    \item How can we use naturally occurring variation to identify causal effects in labor markets, and what do these effects tell us about the functioning of labor market institutions?
\end{itemize}

\textbf{In-Scope:}
\begin{itemize}[nosep]
    \item Fundamental Question
    \item Minimum Wage and Employment
    \item Immigration and Labor Markets
    \item Returns to Education
    \item Wage Inequality and Firm Effects
\end{itemize}

\textbf{Out-of-Scope (delegiert an andere Appendices):}
\begin{itemize}[nosep]
    \item Glossary $\rightarrow$ Appendix G
\end{itemize}

\textbf{Constraints (Anwendungsgrenzen):}
\begin{itemize}[nosep]
    \item [TODO: Define when this appendix does NOT apply]
    \item [TODO: Define required prerequisites]
    \item [TODO: Define known limitations]
\end{itemize}

\textbf{Lieferobjekte (Deliverables):}
\begin{itemize}[nosep]
    \item 1 Table(s)
    \item 6 Equation(s)
    \item Worked Example(s)
\end{itemize}

\end{tcolorbox}


\section{Fundamental Question}

\begin{tcolorbox}[colback=red!5!white, colframe=red!75!black, title=Central Question]
\textbf{How can we use naturally occurring variation to identify causal effects in labor markets, and what do these effects tell us about the functioning of labor market institutions?}
\end{tcolorbox}

Card's answer: By exploiting \textbf{natural experiments}---policy changes, geographic variation, and institutional discontinuities---we can identify causal effects that challenge conventional economic wisdom and inform evidence-based policy.

% ============================================================================
% THEORETICAL FOUNDATIONS
% ============================================================================
\section{Theoretical Foundations: Design-Based Research}

\subsection{The Natural Experiment Revolution}

Card's methodological contribution (Nobel Lecture 2022) emphasizes:

\begin{definition}[Design-Based Research]
An approach to causal inference that:
\begin{enumerate}
    \item Exploits naturally occurring variation (policy changes, geographic borders, cohort differences)
    \item Focuses on \textbf{identification} before estimation
    \item Uses transparent assumptions that can be tested
    \item Prioritizes internal validity with clear scope conditions
\end{enumerate}
\end{definition}

\begin{axiom}[DC-1: Natural Experiment Validity]
\label{ax:dc1}
A natural experiment provides valid causal estimates when:
\begin{equation}
E[Y_0 | D=1] = E[Y_0 | D=0]
\end{equation}
i.e., the treatment assignment is ``as good as random'' conditional on observables.
\textbf{EBF Connection:} Foundation for METHOD-EXPERIMENT causal identification.
\end{axiom}

% ============================================================================
% MINIMUM WAGE
% ============================================================================
\section{Minimum Wage and Employment}

\subsection{The Card-Krueger Revolution}

Card \& Krueger (1994) challenged the textbook prediction that minimum wages reduce employment:

\begin{tcolorbox}[colback=blue!5!white, colframe=blue!75!black]
\textbf{The New Jersey Natural Experiment:}
\begin{itemize}
    \item New Jersey raised minimum wage from \$4.25 to \$5.05 in 1992
    \item Pennsylvania (control) kept minimum wage at \$4.25
    \item Fast-food employment \textbf{increased} in NJ relative to PA
    \item Difference-in-differences estimate: +2.76 FTE employees per restaurant
\end{itemize}
\end{tcolorbox}

\begin{axiom}[DC-2: Monopsony in Labor Markets]
\label{ax:dc2}
The near-zero employment effect of minimum wages suggests labor market monopsony:
\begin{equation}
\frac{\partial L}{\partial w_{min}} \approx 0 \quad \text{when } w_{min} < w^{monopsony}
\end{equation}
Firms have wage-setting power, so moderate minimum wage increases transfer rents without reducing employment.
\textbf{EBF Connection:} Challenges simple supply-demand models in DOMAIN-LABOR.
\end{axiom}

\subsection{Policy Implications}

Card \& Krueger (1995) ``Myth and Measurement'' synthesizes the evidence:
\begin{itemize}
    \item Small minimum wage increases have minimal employment effects
    \item Larger increases may have negative effects at the margin
    \item Effects vary by labor market tightness and firm characteristics
    \item Distributional effects favor low-wage workers
\end{itemize}

% ============================================================================
% IMMIGRATION
% ============================================================================
\section{Immigration and Labor Markets}

\subsection{The Mariel Boatlift}

Card (1990) exploited the sudden influx of 125,000 Cuban immigrants to Miami:

\begin{tcolorbox}[colback=green!5!white, colframe=green!75!black]
\textbf{Mariel Boatlift Natural Experiment:}
\begin{itemize}
    \item 7\% increase in Miami labor force in 1980
    \item Comparison cities: Atlanta, Houston, Los Angeles, Tampa
    \item \textbf{No significant effect} on wages or unemployment of natives
    \item Even for low-skilled workers most similar to Mariel immigrants
\end{itemize}
\end{tcolorbox}

\begin{axiom}[DC-3: Labor Market Absorption]
\label{ax:dc3}
Labor markets can absorb large immigrant inflows through:
\begin{enumerate}
    \item Demand adjustments (industry composition shifts)
    \item Native out-migration to other regions
    \item Capital deepening in response to labor supply
\end{enumerate}
\begin{equation}
\frac{\partial w_{native}}{\partial L_{immigrant}} \approx -0.1 \text{ to } -0.2
\end{equation}
\textbf{EBF Connection:} Informs CONTEXT-DEMOGRAPHIC parameters for policy.
\end{axiom}

\subsection{Places vs. People: Reconciling the Literature}

\citet{PAP-dustmann2025immigration} provide a unified framework that reconciles the seemingly contradictory findings between regional (Card-style) and worker-level (Borjas-style) approaches:

\begin{tcolorbox}[colback=blue!5!white, colframe=blue!75!black]
\textbf{Key Insight: Both Approaches Are Valid But Answer Different Questions}
\begin{itemize}
    \item \textbf{Regional approaches} (Card, 1990): Measure effects on \textit{places}---including displaced natives who leave
    \item \textbf{Worker-level approaches} (Borjas, 2003): Measure effects on \textit{people}---natives who remain in competition
\end{itemize}
\end{tcolorbox}

\textbf{Decomposition of Regional Employment Effects:}
\begin{equation}
\frac{dL^N_r}{dL^I_r} = \underbrace{\text{Displacement}}_{\approx 0.14} + \underbrace{\text{Crowding-out}}_{\approx 0.02} + \underbrace{\text{Relocation}}_{\approx 0.71} = -0.87
\end{equation}

\textbf{Key Parameters from German Data:}
\begin{itemize}
    \item Inverse labor demand elasticity: $\phi^{-1} = -1.95$
    \item Place-level labor supply elasticity: $\eta^P = 4.64$
    \item Worker-level labor supply elasticity: $\eta^E = 3.68$
    \item Pure wage effect: $-0.188$
\end{itemize}

\textbf{Policy Implication:} Neither approach is ``correct''---both provide useful estimates for different policy questions. Regional effects inform area-level policies; worker-level effects inform individual welfare assessments.

% ============================================================================
% RETURNS TO EDUCATION
% ============================================================================
\section{Returns to Education}

\subsection{The Identification Challenge}

Card (1999, 2001) addresses the endogeneity of schooling:

\begin{itemize}
    \item \textbf{Ability bias}: More able individuals get more education
    \item \textbf{Measurement error}: Self-reported education is noisy
    \item \textbf{Selection}: Who chooses to continue schooling?
\end{itemize}

\subsection{Instrumental Variable Solutions}

\begin{axiom}[DC-4: IV Estimates of Returns]
\label{ax:dc4}
Using instruments (college proximity, compulsory schooling laws, twins), the causal return to education is:
\begin{equation}
\frac{\partial \ln(wage)}{\partial years\_education} \approx 0.08-0.12
\end{equation}
IV estimates are often \textbf{higher} than OLS, suggesting ability bias is offset by measurement error attenuation.
\textbf{EBF Connection:} Provides calibration for DOMAIN-EDUCATION human capital models.
\end{axiom}

% ============================================================================
% WAGE INEQUALITY
% ============================================================================
\section{Wage Inequality and Firm Effects}

\subsection{The Role of Firms}

Card, Heining \& Kline (2013) decompose wage inequality using matched employer-employee data:

\begin{equation}
\ln(w_{ijt}) = \alpha_i + \psi_j + X'_{it}\beta + \varepsilon_{ijt}
\end{equation}

where $\alpha_i$ is worker fixed effect, $\psi_j$ is firm fixed effect.

\begin{axiom}[DC-5: Firm Wage Premiums]
\label{ax:dc5}
Firm effects explain 20-25\% of wage variance:
\begin{itemize}
    \item High-wage firms pay \textbf{all} workers more (not just high-skill)
    \item Sorting of high-skill workers to high-paying firms increased over time
    \item Rising inequality = rising firm dispersion + increased sorting
\end{itemize}
\textbf{EBF Connection:} Extends DOMAIN-LABOR with firm-level heterogeneity.
\end{axiom}

\subsection{Peer Effects and Fairness}

Card, Mas, Moretti \& Saez (2012) show that relative pay matters:

\begin{tcolorbox}[colback=yellow!10!white, colframe=yellow!75!black]
\textbf{Peer Salary Experiment:}
\begin{itemize}
    \item UC employees given access to colleagues' salaries
    \item Those earning \textbf{below} median: job satisfaction drops 15\%
    \item Those earning \textbf{above} median: no change in satisfaction
    \item Asymmetric response consistent with loss aversion
\end{itemize}
\end{tcolorbox}

\begin{axiom}[DC-6: Social Comparison in Wages]
\label{ax:dc6}
Utility from wages includes a social comparison component:
\begin{equation}
U(w_i) = u(w_i) + \lambda \cdot \mathbb{1}_{w_i < w_{reference}} \cdot (w_{reference} - w_i)
\end{equation}
where $\lambda > 0$ captures loss aversion to being below reference.
\textbf{EBF Connection:} Links to CORE-WHAT fairness dimension and loss aversion.
\end{axiom}

% ============================================================================
% ACTIVE LABOR MARKET POLICY
% ============================================================================
\section{Active Labor Market Policy}

\subsection{Meta-Analytic Evidence}

Card, Kluve \& Weber (2010, 2018) synthesize evidence from 200+ program evaluations:

\begin{tcolorbox}[colback=blue!5!white, colframe=blue!75!black, title=What Works in Labor Market Policy?]
\begin{tabular}{lcc}
\textbf{Program Type} & \textbf{Short-run} & \textbf{Long-run} \\
\hline
Job search assistance & ++ & + \\
Training programs & 0/+ & ++ \\
Private sector subsidies & + & + \\
Public employment & 0 & 0 \\
\end{tabular}
\begin{itemize}
    \item Training programs: negative short-run (lock-in), positive long-run
    \item Public employment programs: consistently ineffective
    \item Effects larger for women, long-term unemployed
\end{itemize}
\end{tcolorbox}

% ============================================================================
% EBF INTEGRATION
% ============================================================================
\section{EBF Integration}

\subsection{10C Mapping}

\begin{table}[h]
\centering
\begin{tabular}{lll}
\textbf{10C Element} & \textbf{Card Contribution} & \textbf{Parameter} \\
\hline
CORE-WHO (AAA) & Worker/firm heterogeneity & Firm effects $\psi_j$ \\
CORE-WHAT (C) & Fairness in wages & Social comparison $\lambda$ \\
CORE-HOW (B) & Labor market complementarities & Skill-firm matching \\
CORE-WHEN (V) & Policy timing effects & Short vs. long-run \\
CORE-WHERE (BBB) & Labor market parameters & Returns, elasticities \\
CORE-AWARE (AU) & Wage transparency effects & Peer salary knowledge \\
CORE-READY (AV) & Job search intensity & ALMP effectiveness \\
CORE-STAGE (AW) & Career progression & Education-wage path \\
\end{tabular}
\caption{Card contributions mapped to 10C framework}
\end{table}

\subsection{Key Parameter Values for BBB}

\begin{tcolorbox}[colback=green!5!white, colframe=green!75!black, title=Card Parameters for BBB Repository]
\begin{itemize}
    \item Returns to education: 8-12\% per year (IV estimates)
    \item Minimum wage employment elasticity: $\approx 0$ (for moderate increases)
    \item Immigration wage elasticity: -0.1 to -0.2
    \item Firm wage premium variance share: 20-25\%
    \item Peer salary dissatisfaction: -15\% if below median
    \item Training program long-run effect: +5-10\% earnings
\end{itemize}
\end{tcolorbox}

% ============================================================================
% WORKED EXAMPLE
% ============================================================================
\section{Worked Example: Minimum Wage Policy Analysis}

Consider a city contemplating a minimum wage increase from \$12 to \$15:

\textbf{Card Framework Analysis:}

\begin{enumerate}
    \item \textbf{Identify comparison group}: Neighboring city without increase
    \item \textbf{Estimate baseline}: Current employment in affected industries
    \item \textbf{Apply Card-Krueger elasticity}: $\epsilon \approx 0$ for moderate increase
    \item \textbf{Predict distributional effects}:
    \begin{itemize}
        \item Workers earning \$12-15: wage increase to \$15
        \item Workers earning >\$15: potential spillover (+2-5\%)
        \item Employment: minimal change predicted
    \end{itemize}
    \item \textbf{Monitor with difference-in-differences}: Compare trends post-policy
\end{enumerate}

\textbf{Caveat}: Large increases (>30\%) may show negative employment effects.

% ============================================================================
% SUMMARY
% ============================================================================


%%%%%%%%%%%%%%%%%%%%%%%%%%%%%%%%%%%%%%%%%%%%%%%%%%%%%%%%%%%%%%%%%%%%%%%%%%%%%%%
\section{Key Results}
\label{sec:dc-results}
%%%%%%%%%%%%%%%%%%%%%%%%%%%%%%%%%%%%%%%%%%%%%%%%%%%%%%%%%%%%%%%%%%%%%%%%%%%%%%%

The analysis yields the following key results:

\begin{enumerate}
    \item \textbf{Result 1:} Primary finding with quantitative support
    \item \textbf{Result 2:} Secondary finding with implications
    \item \textbf{Result 3:} Practical application or boundary condition
\end{enumerate}

\section{Summary}

\begin{tcolorbox}[colback=gray!10!white, colframe=gray!75!black, title=Key Takeaways]
\begin{enumerate}
    \item \textbf{Natural Experiments:} Exploit policy variation for causal identification
    \item \textbf{Minimum Wage:} Near-zero employment effects for moderate increases
    \item \textbf{Immigration:} Labor markets absorb large inflows with small wage effects
    \item \textbf{Returns to Education:} 8-12\% per year, robust to endogeneity corrections
    \item \textbf{Firm Effects:} 20-25\% of wage variance, increasing sorting over time
    \item \textbf{ALMP:} Training works long-run; public employment doesn't work
\end{enumerate}
\end{tcolorbox}

% ============================================================================
% GLOSSARY
% ============================================================================
\section{Glossary}

See Appendix G (Glossary) for standard EBF terms. Card-specific terms:

\begin{description}
    \item[Design-Based Research] Causal inference prioritizing identification through research design
    \item[Difference-in-Differences] Method comparing treatment-control differences before/after intervention
    \item[Firm Fixed Effect] Time-invariant wage premium associated with a specific employer
    \item[Monopsony] Labor market with wage-setting power by employers
    \item[Natural Experiment] Naturally occurring variation that mimics random assignment
\end{description}

% ============================================================================
% CRITICAL FOUNDATIONS
% ============================================================================
\section{Critical Foundations}

\begin{enumerate}
    \item Card, D. \& Krueger, A. (1994). Minimum Wages and Employment. \textit{AER}. [Revolutionary]
    \item Card, D. (1990). The Impact of the Mariel Boatlift. \textit{ILRR}. [Natural Experiment Classic]
    \item Card, D. (1999). The Causal Effect of Education on Earnings. \textit{Handbook}. [Survey]
    \item Card, D. et al. (2013). Workplace Heterogeneity and Wage Inequality. \textit{QJE}. [Firm Effects]
    \item Card, D. et al. (2018). What Works? ALMP Meta-Analysis. \textit{JEEA}. [Policy]
\end{enumerate}

% ============================================================================
% OPEN ISSUES
% ============================================================================
\section{Open Issues}

\begin{enumerate}
    \item \textbf{External validity:} Do natural experiment results generalize to other contexts?
    \item \textbf{Minimum wage thresholds:} At what level do negative employment effects emerge?
    \item \textbf{Immigration long-run:} What are the dynamic effects beyond short-run adjustment?
    \item \textbf{Firm effects mechanisms:} Why do some firms pay more? Rent-sharing vs. sorting?
\end{enumerate}

% ============================================================================
% REFERENCES
% ============================================================================
\section{References}

\nocite{bcm_master}

For complete bibliography, see Master Bibliography (\texttt{bibliography/master\_references.bib}).

Key Card references integrated in this appendix:
\begin{itemize}[nosep]
    \item card1994minimum, card1990mariel, card1995myth
    \item card1999causal, card2001estimating, card2012inequality
    \item card2012peer, card2015active, card2018active2
    \item card2022nobel
\end{itemize}

\end{document}
